% Uma alternativa que visa superar essas limitações são os modelos baseados em dados, que utilizam técnicas para
% identificar e explorar relações estatísticas entre os dados de entrada e saída. Com a aprendizagem dos padrões
% presentes nos dados históricos, é possível mapear as dinâmicas subjacentes do sistema e utilizá-las para prever
% valores futuros. Assim, esses modelos são capazes de realizar previsões precisas com base em séries temporais,
% capturando tendências e comportamentos complexos que podem ser difíceis de modelar diretamente por meio de abordagens
% tradicionais. Diferentemente das mesmas, que demandam tempo e dados específicos, os modelos baseados em dados oferecem
% maior agilidade para previsões em tempo real, sendo uma escolha mais adequada para sistemas de alerta precoce e
% mitigação de desastres de cheias \cite{ta}.

% Uma das maiores vantagens de utilizar essa abordagem é o grande desenvolvimento do campo de Análise e Predição de
% Séries Temporais. Com isso, é possível utilizar as técnicas desenvolvidas para a aplicação em problemas com uma
% modelagem semelhante, mas em contextos distintos \cite{hyndman}, como Economia, Biologia, entre outros. Dentro deste
% contexto, uma ampla gama de modelagens são utilizadas para a predição de séries temporais. Isso apresenta uma
% oportunidade de análise e comparação entre as vantagens e desvantagens das mesmas com o intuíto de descobrir a mais
% efetiva para a situação particular \cite{makridakis, morales}.

% Uma abordagem mais clássica e altamente implementada destes modelos vem na forma de modelos estatísticos, como o ARIMA
% (\textit{Autoregressive Integrated Moving Average}) \cite{box}, um modelo que descreve processos estocásticos usando
% uma combinação de autoregressão, isto é, a predição de uma Série Temporal usando somente valores anteriores da mesma,
% a média móvel - ajuste das previsões considerando os erros residuais das observações anteriores, o que permite
% capturar flutuações - e, por último, a diferenciação dos valores consecutivos da série, reduzindo o impacto de padrões
% de longo prazo e tendências. Além disso, um elemento de sazionalidade, ou seja, padrões repetidos em com uma certa
% frequência fixa, como padrões mensais ou trimestrais, podem ser compreendidos através de uma extensão do ARIMA, o
% SARIMA. Historicamente, esses modelos têm sido usados em predições em contextos semelhantes como em Atashi et al.
% (2022) \cite{atashi} que comparou o desempenho do SARIMA com os modelos LSTM e \textit{Random Forest} para predição do
% nível da água no Rio Vermelho do Norte, localizado na fronteira entre os Estados Unidos e o Canadá, e Güldal et al.
% (2009) \cite{gudal} que utilizou o ARMA para a predição do Lago Eğirdir na Túrquia, oferecendo uma base para
% comparações com métodos modernos.

% Esses modelos, eficazes e de baixo custo computacional, apresentam tanto limitações quanto vantagens quando comparados
% com modelos mais recentes desenvolvidos nas áreas de Aprendizado de Máquina (\textit{Machine Learning}) e Aprendizado
% Profundo (\textit{Deep Learning}) \cite{kontopoulou}. No âmbito de \textit{Machine Learning}, modelos baseados em
% \textit{Gradient Boosting} têm se mostrado promissores em situações de desastres, onde a rapidez e precisão nas
% previsões são essenciais para mitigar impactos \cite{albahri}. Um exemplo amplamente utilizado é o XGBoost
% (\textit{eXtreme Gradient Boosting}) \cite{chen_xgboost}. O mesmo utiliza árvores de decisão em série, ajustando
% constantemente os erros das previsões anteriores. Além disso, ele incorpora regularização para reduzir o risco de
% sobreajuste (\textit{overfitting}) e otimizações específicas para lidar com grandes volumes de dados. Uma
% implementação semelhante, visando reduzir o tempo computacional, é dada pelo LightGBM \cite{lightgbm}. Ela mantém as
% vantagens do XGBoost, como otimização para dados esparsos, treinamento paralelo, múltiplas funções de perda e
% regularização, mas altera o método de construção das árvores de decisão, mostrando-se uma boa alternativa para obter
% resultados semelhantes de forma mais eficaz, algo imprescindível no contexto de uma crise climática como a ocorrida.

% Por outro lado, em Aprendizado Profundo, modelos baseados em Redes Neurais Recorrentes (RNNs), como o \textit{Long
% Short Term Memory} (LSTM), são amplamente utilizados em situações semelhantes, como demonstrado por Borwarnginn et al.
% \cite{borwarnginn}. O LSTM foi projetado especificamente para lidar com dependências de longo prazo em séries
% temporais, garantindo que informações relevantes de eventos passados sejam preservadas, ao mesmo tempo que mantém o
% contexto de eventos mais recentes. Essa análise é particularmente útil em problemas onde padrões temporais complexos
% e, geralmente, não lineares precisam ser modelados.